\chapter{Tests, résultats et étude de cas}
\label{chap:tests}

\callout{En bref}{Ce chapitre présente la stratégie de validation du système,
les résultats des tests automatisés, puis une étude de cas approfondie sur un
jeu de 17 769 matchs internationaux de football, avant une discussion des
apports et des limites.}

\section{Stratégie de tests}
La qualité du système repose sur une suite de \textbf{22 tests automatisés}
(exécutés avec \texttt{pytest}), couvrant à la fois le moteur déterministe et
la logique d'orchestration. Deux principes ont guidé cette stratégie :
\begin{itemize}
  \item \textbf{tester sans LLM} : le cœur déterministe est validé
        indépendamment de tout fournisseur de modèle, ce qui rend les tests
        rapides et reproductibles ;
  \item \textbf{simuler les défaillances d'agents} : un client LLM
        « fictif » permet de vérifier que les filets de sécurité produisent
        des livrables complets même lorsque les agents n'appellent aucun outil.
\end{itemize}

Le tableau~\ref{tab:tests} récapitule la couverture par module.

\begin{table}[ht]
\centering\small
\caption{Couverture des tests automatisés.}
\label{tab:tests}
\begin{tabularx}{\linewidth}{@{}Yc@{}}
\toprule
\textbf{Domaine testé} & \textbf{Nb. de tests} \\
\midrule
Profilage des rôles de colonnes & 1 \\
Suggestion et calcul des KPIs & 1 \\
Suggestion et construction des graphiques & 1 \\
Assemblage du tableau de bord HTML & 1 \\
Permissions du serveur MCP (refus hors périmètre) & 1 \\
Vérification des chiffres (\texttt{validation.py}) & 7 \\
Relance avec incitation (\textit{retry-with-nudge}) & 2 \\
Détection géographique \& carte choroplèthe & 3 \\
Détection d'anomalies & 2 \\
Comparaison de segments & 2 \\
Extraction numérique robuste & 1 \\
\midrule
\textbf{Total} & \textbf{22} \\
\bottomrule
\end{tabularx}
\end{table}

L'ensemble des 22 tests est au vert. Au-delà des tests unitaires, des
\textbf{exécutions de bout en bout} sur des jeux de données réels valident
l'intégration complète (du téléversement aux artefacts produits).

\section{Étude de cas : matchs internationaux de football}

\subsection{Présentation du jeu de données}
Pour valider l'approche généraliste, nous avons soumis au système un jeu de
données de \textbf{17 769 matchs internationaux de football} disputés entre
\textbf{1872 et 2024}. Aucune configuration spécifique n'a été fournie : le
système a découvert seul la structure des données. Le profilage automatique a
identifié :
\begin{itemize}
  \item 2 \textbf{mesures} : buts à domicile (\textit{Home Goals}) et à
        l'extérieur (\textit{Away Goals}) ;
  \item 2 colonnes \textbf{géographiques} : équipe à domicile et à
        l'extérieur (reconnues comme des pays) ;
  \item 1 colonne \textbf{date}, et plusieurs \textbf{dimensions}
        (tournoi, ville, stade neutre ou non) ;
  \item 90 tournois, 214 équipes à domicile et 218 à l'extérieur.
\end{itemize}

La figure~\ref{fig:source} montre la section « source des données » du rapport
généré, incluant le schéma détecté et le résultat de la validation.

\begin{figure}[ht]
\centering
\includegraphics[width=0.92\linewidth]{fig_rapport_source.png}
\caption{Section « source des données » : schéma détecté automatiquement et validation.}
\label{fig:source}
\end{figure}

\subsection{Résultats : indicateurs et tableau de bord}
Le système a produit automatiquement un tableau de bord complet et un rapport.
Parmi les indicateurs calculés :
\begin{itemize}
  \item un total de \textbf{30 605 buts à domicile} sur l'ensemble des matchs ;
  \item une moyenne de \textbf{1,72} but à domicile contre \textbf{1,17} à
        l'extérieur ;
  \item un maximum de \textbf{31 buts} marqués en un seul match.
\end{itemize}

Le tableau~\ref{tab:resultats} récapitule les principaux indicateurs produits
automatiquement par le système pour ce jeu de données.

\begin{table}[ht]
\centering\small
\caption{Indicateurs clés calculés pour l'étude de cas.}
\label{tab:resultats}
\begin{tabularx}{\linewidth}{@{}Yr@{}}
\toprule
\textbf{Indicateur} & \textbf{Valeur} \\
\midrule
Nombre total de matchs & 17\,769 \\
Total des buts à domicile & 30\,605 \\
Moyenne de buts à domicile & 1,72 \\
Moyenne de buts à l'extérieur & 1,17 \\
Maximum de buts en un match & 31 \\
Tournois distincts & 90 \\
Équipes à domicile distinctes & 214 \\
Équipes à l'extérieur distinctes & 218 \\
\bottomrule
\end{tabularx}
\end{table}

La \textbf{carte choroplèthe} (figure~\ref{fig:carte}, chapitre précédent)
révèle d'un coup d'œil la domination du \textbf{Brésil} en buts cumulés à
domicile, suivi des grandes nations européennes et sud-américaines. Les
\textbf{séries temporelles} (figures~\ref{fig:serieH} et~\ref{fig:serieA})
montrent une nette intensification du nombre de buts après l'an 2000, reflet de
la multiplication des rencontres internationales.

\subsection{Résultats : anomalies détectées}
La détection d'anomalies a automatiquement isolé des matchs au score
exceptionnel. Le plus extrême correspond au match historique
\textbf{Australie 31--0 Samoa américaines} (2001), record du monde de buts
inscrits en un match international, correctement identifié comme valeur
aberrante (score~z très élevé). Côté extérieur, le match \textbf{Maldives
0--17 Iran} (1997) ressort également. Ce résultat illustre la capacité du
système à faire remonter, sans aucune connaissance métier préalable, des
événements réellement remarquables.

\subsection{Résultats : avantage du terrain}
La comparaison de segments a mis en évidence un phénomène bien connu des
analystes du football : l'\textbf{avantage du terrain}. En segmentant selon le
caractère « à domicile ou non » du stade, le système a mesuré une moyenne de
buts supérieure pour l'équipe jouant chez elle (1,78 contre 1,48), confirmant
quantitativement cet effet — là encore, de manière entièrement automatique.

\subsection{Véracité du rapport}
Enfin, le mécanisme de vérification a confirmé que \textbf{l'intégralité des
chiffres} cités dans les conclusions du rapport correspondait aux indicateurs
réellement calculés (verdict « vérifié », figure~\ref{fig:findings} du
chapitre précédent), démontrant l'efficacité du garde-fou anti-hallucination.

\section{Discussion}

\subsection{Apports}
L'étude de cas confirme les promesses de l'approche :
\begin{itemize}
  \item \textbf{généralité} : un jeu de données sportif, jamais anticipé, est
        analysé sans aucune configuration ;
  \item \textbf{pertinence} : les indicateurs, la carte et les insights
        produits sont immédiatement exploitables ;
  \item \textbf{fiabilité} : les chiffres sont exacts et vérifiés, et les
        livrables toujours complets grâce aux filets de sécurité ;
  \item \textbf{rapidité} : l'analyse complète est produite en quelques
        secondes, contre des heures de travail manuel.
\end{itemize}

\subsection{Limites}
Plusieurs limites subsistent :
\begin{itemize}
  \item la détection géographique repose sur des noms de pays en anglais ;
        certaines entités infranationales (Écosse, Pays de Galles) ne sont pas
        cartographiées par la projection mondiale ;
  \item la qualité de la narration dépend du modèle de langage employé ;
  \item le système traite un fichier à la fois et ne gère pas encore les
        jointures entre sources ;
  \item l'analyse reste descriptive (pas de modélisation prédictive).
\end{itemize}
Ces limites ouvrent directement sur les perspectives présentées en conclusion.

\callout{Synthèse du chapitre}{Les 22 tests valident le cœur du système, et
l'étude de cas sur 17 769 matchs démontre une analyse généraliste, pertinente
et fiable : carte mondiale, détection du match record 31--0, mise en évidence
de l'avantage du terrain, et chiffres intégralement vérifiés.}
